\documentclass[10pt,a4paper]{amsart}
\usepackage[margin=19mm]{geometry}
\usepackage{amsmath,amssymb,mathtools}
\usepackage[scr=rsfs,cal=euler]{mathalfa}
\usepackage{xfrac}
\usepackage[unicode,hidelinks,bookmarks=false]{hyperref}
\newcommand{\ie}{i.e.\ }
\newcommand{\cf}{cf.\ }
\newcommand{\resp}{respectively\ }
\newcommand{\Subsection}{\subsection}
\providecommand{\nobreakdash}{\nobreak\hbox{-}}
\setlength{\emergencystretch}{2em}
\multlinegap0pt
\newtheorem{comment}{Comment}
\usepackage{amssymb}
\usepackage{tikz}
\usepackage{cancel}
\usepackage{xcolor}
\newtheorem{theorem}{Theorem}
\newcommand{\shalmalirmk}[1]{\textcolor{red}{\textsf{[Shalmali: #1]}}}
\title{Nonlinear elliptic Dirichlet boundary value problems on time scales}
\author{Shalmali Bandyopadhyay, F. Ay\c{c}a \c{C}etinkaya, Tom Cuchta}
\date{Source: arXiv:2602.10335v1 (CC BY 4.0)}
\begin{document}
\maketitle

\noindent\emph{Source and licence.} Nonlinear elliptic Dirichlet boundary value problems on time scales. Version arXiv:2602.10335v1 (CC BY 4.0), distributed by arXiv under the Creative Commons Attribution 4.0 International licence, http://creativecommons.org/licenses/by/4.0/, as recorded in the arXiv licence metadata for this exact version and shown on the abstract page https://arxiv.org/abs/2602.10335v1. Full licence notice: \texttt{licenses/arxiv-2602.10335-CC-BY-4.0.txt}.\par\smallskip

\noindent\textit{Editorial scope.} Counted unit: the article's theorem thm:A-inverse (source0.tex lines 408--414). The article's complete original proof of it is delivered with it (source0.tex lines 416--467, one \texttt{proof} environment of 52 lines). No mathematics is written, repaired or re-derived: the statement and the proof are the article's own lines, in the article's own notation, and no retained line is substituted or re-flowed.  Retained uncounted as the source-local context the proof needs: lines 240--242; lines 391--396.  Nothing was omitted from any retained span, so the package carries no omission record.  No retained line contains a citation, so the wrapper carries no bibliography and no bibliography entry.  No retained line contains a figure environment, an image-inclusion command or drawing code, so no image is delivered and the package declares no asset dependency.  Editorial formatting only, in addition: an AMS \texttt{amsart} preamble that restates the article's own declarations and macros used by the retained lines, namely the environments \texttt{theorem} on separate counters, together with the standard packages those lines need.  The retained environments print in this document as Theorem 1 in order of appearance, whereas the article prints the same items with the numbering of its own full document.  The complete native source of the article as distributed in the arXiv e-print for this exact version is bundled unchanged as \texttt{sources/194387/source0.tex}.
\begin{equation}\label{eq:nabla-from-delta}
f^\nabla(t) = f^\Delta(\rho(t)) \quad \text{for all } t \in (a,b]\cap\mathbb{T};
\end{equation}

\begin{equation}\label{eq:Green-1D}
G(t,s) = \frac{1}{b-a} \begin{cases}
(t-a)(b-s) & \text{if } t \leq s,\\
(s-a)(b-t) & \text{if } t \geq s.
\end{cases}
\end{equation}

\begin{theorem}\label{thm:A-inverse}
Let $G$ be defined as in \eqref{eq:Green-1D}. Then the inverse operator $A^{-1}\colon H \to D$ is
\begin{equation}\label{eq:A-inverse-rep}
\left(A^{-1}u\right)(t) = \int_a^b G(t,s) u(s) \Delta s.
\end{equation}
Moreover, $A^{-1}$ is symmetric and compact on $H$.
\end{theorem}

\begin{proof}
Let $u \in H$ and define $y(t) := \displaystyle\int_a^b G(t,s) u(s) \, \Delta s$. From \eqref{eq:Green-1D}, $G(a,s)=G(b,s)=0$ for all $s$, hence $y(a)=y(b)=0$. For $t\in(a,b)\cap\mathbb{T}$,
$$y(t) = \int_a^t G(t,s)u(s)\,\Delta s + \int_t^b G(t,s)u(s)\,\Delta s = \frac{b-t}{b-a}\int_a^t (s-a)u(s)\,\Delta s + \frac{t-a}{b-a}\int_t^b (b-s)u(s)\,\Delta s.$$
Define $I_1(t)=\displaystyle\int_a^t (s-a)u(s)\,\Delta s$ and $I_2(t)=\displaystyle\int_t^b (b-s)u(s)\,\Delta s$, so that $y(t)=\dfrac{1}{b-a}\Big[(b-t)I_1(t)+(t-a)I_2(t)\Big]$. Then $I_1^\Delta(t)=(t-a)u(t)$ and $I_2^\Delta(t)=-(b-t)u(t)$, hence after the product rule and simplifying,
\begin{equation}\label{eq:y-delta-correct}
(b-a)y^\Delta(t)=I_2(t)-I_1(t)-(b-a)\mu(t)u(t).
\end{equation}
Applying \eqref{eq:nabla-from-delta}, \eqref{eq:y-delta-correct} becomes
$$(b-a)y^\nabla(t) = I_2(\rho(t))-I_1(\rho(t))-(b-a)\mu(\rho(t))u(\rho(t)).$$
To complete the verification, we compute $y^{\nabla\Delta}$ directly. By \eqref{eq:nabla-from-delta}, $I_1^\nabla(t) = I_1^\Delta(\rho(t)) = (\rho(t)-a)u(\rho(t))$ and $I_2^\nabla(t) = I_2^\Delta(\rho(t)) = -(b-\rho(t))u(\rho(t))$. By the product rule,
\begin{align*}
[(b-t)I_1(t)]^\nabla &= -I_1(t) + (b-\rho(t))(\rho(t)-a)u(\rho(t)), \\
[(t-a)I_2(t)]^\nabla &= I_2(t) - (\rho(t)-a)(b-\rho(t))u(\rho(t)).
\end{align*}
Adding these expressions reveals $(b-a)y^\nabla(t) = I_2(t) - I_1(t)$. Now we compute
\[y^{\nabla\Delta}(t) = \dfrac{I_2^\Delta(t) - I_1^\Delta(t)}{b-a}=\frac{-(b-t)u(t) - (t-a)u(t)}{b-a}=-u(t).\]
Hence $-y^{\nabla\Delta}(t) = u(t)$ for $t \in (a,b)$, and thus $y = A^{-1}u$.

\begin{comment}
To complete the verification, we compute $y^{\nabla\Delta}$ directly. By \eqref{eq:nabla-from-delta}, $y^\nabla(t) = y^\Delta(\rho(t))$, so from \eqref{eq:y-delta-correct},
\begin{equation}\label{eq:y-nabla-formula}
(b-a)y^\nabla(t) = I_2(\rho(t)) - I_1(\rho(t)) - (b-a)\mu(\rho(t))u(\rho(t)).
\end{equation}
Since $\rho(\sigma(t)) = t$ for $t \in (a,b) \cap \mathbb{T}$, 
%\shalmalirmk{this is not true}
we have
\[
(b-a)y^\nabla(\sigma(t)) = I_2(t) - I_1(t) - (b-a)\mu(t)u(t).
\]
Computing the difference and using the fundamental theorem of calculus on time scales,
\[
I_1(t) - I_1(\rho(t)) = \nu(t)(\rho(t)-a)u(\rho(t)), \quad I_2(t) - I_2(\rho(t)) = -\nu(t)(b-\rho(t))u(\rho(t)).
\]
Substituting and simplifying,
\[
y^\nabla(\sigma(t)) - y^\nabla(t) = -\nu(t)u(\rho(t)) - \mu(t)u(t) + \mu(\rho(t))u(\rho(t)).
\]
Using the identity $\mu(\rho(t)) = \nu(t)$, division by $\mu(t)$ yields $y^{\nabla\Delta}(t) = -u(t)$. Hence $-y^{\nabla\Delta}(t) = u(t)$ for $t \in (a,b)$, and thus $y = A^{-1}u$.


To complete the verification, we use the {\color{violet} (Tom: see above) defining property} of the Green's function: $G$ satisfies the homogeneous equation $-G^{\nabla_t\Delta_t} = 0$ in the regions $t < s$ and $t > s$, is continuous at $t = s$, and has a jump discontinuity in $G^{\nabla}$ at $t = s$ of magnitude $-1$, i.e.
$$G^\nabla(s^+, s) - G^\nabla(s^-, s) = -\frac{s-a}{b-a} - \frac{b-s}{b-a} = -1.$$
{\color{violet} Tom: it's not clear to me right here why this follows from (3.8)} Hence $-y^{\nabla\Delta}(t)=u(t)$ for $t\in(a,b)$, and thus $y=A^{-1}u$.
\end{comment}
\smallskip

For compactness, let $M=\max\limits_{(t,s)\in\left([a,b]\cap\mathbb{T}\right)^2}|G(t,s)|$. By the Cauchy--Schwarz inequality,
$$|(A^{-1}u)(t)|^2 \le \left(\int_a^b |G(t,s)|^2\,\Delta s\right)\|u\|^2 \le M^2(b-a)\|u\|^2.$$
Hence $\|A^{-1}u\|_\infty\le M\sqrt{b-a}\,\|u\|$. Moreover, for $t_1,t_2\in[a,b]$,
$$\Big|(A^{-1}u)(t_1)-(A^{-1}u)(t_2)\Big| \le \sqrt{b-a}\,\|u\| \left(\int_a^b |G(t_1,s)-G(t_2,s)|^2\,\Delta s\right)^{1/2}.$$
Since $G$ is uniformly continuous on $\left([a,b]\cap\mathbb{T}\right)^2$, the right-hand side tends to $0$ uniformly as $t_1\to t_2$. Thus $A^{-1}$ maps bounded sets of $H$ into equicontinuous, uniformly bounded sets in $C([a,b]\cap\mathbb{T})$. By the Arzel\`a--Ascoli theorem, $A^{-1}$ is compact on $H$. Since $G(t,s)=G(s,t)$ by \eqref{eq:Green-1D}, we have $\langle A^{-1}u, v\rangle = \langle u, A^{-1}v\rangle$, so $A^{-1}$ is symmetric.
\end{proof}

\end{document}
